\section{Formal verification and reproducibility}\label{sec:verification}

\subsection{What the formal endpoints assert}

The Lean definition \path{LogarithmExtension.EventualLowerBound}
is exactly the predicate \(\ELB\) in
Section~\ref{sec:preliminaries}; it quantifies over every integer
numerator and every sufficiently large natural denominator.
The main declarations include
\begin{Verbatim}[fontsize=\small]
logThree_eventual_lower_bound
logThree_irrationalityExponent_eq_two
rationalLog_eventual_lower_bound
rationalLog_irrationalityExponent_eq_two
algebraicLog_irrationalityExponent_bounds
CompatibleLog.compatible_embedding_eventual_lower_bound
CompatibleLog.compatible_embedding_irrationalityExponent_bounds
CompatibleLog.all_embeddings_irrationalityExponent_eq_two
\end{Verbatim}
All are in the namespace \texttt{LogarithmExtension}.
The unbundled compatible-embedding theorem takes only the number field,
\(\gamma,\xi,x\), their stated nonzero conditions, a nonempty finite
set of embeddings, and the exponential identities
\eqref{eq:compatibility}. It has no interpolation, nonvanishing,
analytic-bound, or irrationality premise.
The bundled data structure \path{CompatibleLog.Base} contains these
same mathematical data, not a hidden proof of the conclusion.

The names \path{LogThreeTarget}, \path{RationalLogarithmTarget},
and \path{AlgebraicLogarithmTarget} are proposition definitions.
They are not axioms: the corresponding \texttt{target} declarations
prove them.

The exponent definition reused from \texttt{OAI.PiExponent} is a
real supremum of the positive exponents admitting infinitely many
good reduced rational approximants. For every finite-bound endpoint,
the code first proves that this set is bounded above, and that it
contains \(2\). The resulting supremum therefore agrees with the
usual extended-real definition in this paper.
The Liouville counterexample is expressed as absence of any finite
eventual bound, rather than as a possibly misleading real supremum.

\subsection{Correspondence of arguments and source files}

All extension paths below are relative to
\texttt{LogarithmExtension/}. The earlier dedicated
\texttt{LogThree}, \texttt{RationalLog}, and \texttt{AlgebraicLog}
pipelines remain checked; the \texttt{CompatibleLog} pipeline supplies
the unified theorem proved here.

{\small
\let\path\leanfile
\begin{longtable}{@{}>{\raggedright\arraybackslash}p{.37\textwidth}
                     >{\raggedright\arraybackslash}p{.59\textwidth}@{}}
\toprule
Mathematical component & Principal extension modules\\
\midrule
\endfirsthead
\toprule
Mathematical component & Principal extension modules\\
\midrule
\endhead
\bottomrule
\endfoot
Actual target and finite exponent semantics
& \path{Target.lean}, \path{Endpoint.lean},
\path{GeneralExponent.lean}\\[4pt]
Normalized logarithmic coordinates and rigidity
& \path{NormalizedLogJet.lean}, \path{CoordinateScaling.lean},
\path{CurveRigidity.lean}\\[4pt]
Proposition~\ref{prop:curve}, distinct-center case
& \path{CurveContactFamily.lean},
\path{CurveDerivativeVanishing.lean}, \path{CurveInequality.lean},
\path{CurveGeometry.lean}\\[4pt]
Common-fiber branch
& \path{CommonFiberCurve.lean}, \path{GeometryData.lean}\\[4pt]
Projective degree, local ideals, and exceptional degree
& \path{CurveModelDegrees.lean}, \path{NormalizedJetIdeal.lean},
\path{CompactNormalizedJetIdeal.lean}, \path{LocalContactLength.lean},
\path{ExceptionalDegree.lean}\\[4pt]
Ampleness and global jet generation
& \path{BlowupGeometry.lean}, \path{BlowupMargin.lean},
\path{BlowupAmple.lean}, \path{JetAmpleness.lean},
\path{GlobalJetPackets.lean}, \path{JetSurjectivity.lean}\\[4pt]
Identification with the actual matrix
& \path{BoundedJetMatrix.lean}, \path{MatrixIdentification.lean},
\path{CompatibleLogGeometry.lean}\\[4pt]
Equation~\eqref{eq:row-expansion} and analytic centers
& \path{AnalyticTranslation.lean},
\path{CompatibleLogAnalyticCenters.lean}\\[4pt]
Arithmetic over \(F\), integral norm, and averaging
& \path{NumberFieldMatrix.lean}, \path{NumberFieldNorm.lean},
\path{NumberFieldDeterminant.lean},
\path{NumberFieldNormAveraging.lean},
\path{CompatibleLogMatrixArithmetic.lean}\\[4pt]
Actual summands and their finite sum
& \path{CompatibleLogLiteralAnalyticSummand.lean},
\path{CompatibleLogAnalyticAggregate.lean}\\[4pt]
Incompatible embeddings and full norm bound
& \path{ConjugateAnalyticBound.lean},
\path{CompatibleLogAllEmbeddings.lean}\\[4pt]
Choices and determinant contradiction
& \path{CompatibleLogSuccessiveApproximations.lean},
\path{CompatibleLogAdmissibleParameters.lean},
\path{CompatibleLogDeterminantContradiction.lean}\\[4pt]
Root-of-unity reduction and final theorem
& \path{CompatibleLogBase.lean}, \path{CompatibleLogNonTorsion.lean},
\path{CompatibleLogMain.lean}\\[4pt]
Theorem~\ref{thm:rational}, including the literal \(\log3\)
& \path{LogThreeMain.lean}, \path{RationalLogMain.lean}\\[4pt]
Corollaries~\ref{cor:algebraic}--\ref{cor:quadratic}
& \path{AlgebraicLogMain.lean}, \path{RationalAffine.lean},
\path{RationalPowerLog.lean}, \path{RationalArctangent.lean},
\path{QuadraticNormalizedLog.lean}\\[4pt]
Proposition~\ref{prop:counterexamples}
& \path{LogarithmCounterexamples.lean}\\
\end{longtable}
}

General results reused from the pinned upstream development include
the persistent-component and weighted multiplicity machinery,
projective and local algebra, numerical ampleness and vanishing,
simplex counts, logarithmic row algebra, collision estimates, and
the approximation endpoint lemmas. Theorem \(\mu(\pi)=2\) itself
is not the hypothesis from which the logarithm result is deduced.
The new target-independent geometry data remove the
\(\pi\)-specific approximation hypotheses from the geometric interface.

\subsection{Axiom audit and kernel replay}

The build and transitive axiom audit use the following commands:
\begin{Verbatim}[fontsize=\small]
lake build -q LogarithmExtension.Verification
lake env lean --trust=0 -Ddebug.skipKernelTC=false \
  LogarithmExtension/Verification.lean
\end{Verbatim}
The audit checked \(1326\) extension theorem declarations.
Every transitive axiom was one of
\[
 \texttt{propext},\qquad\texttt{Classical.choice},\qquad
 \texttt{Quot.sound}.
\]
It also audited the reference \(\pi\) endpoint separately.
A source scan found no proof admissions, new axiom declarations,
unsafe or partial declarations, or \texttt{native\_decide} in the
extension. A defined target proposition is not counted as a proof.

The full combined replay was also rerun successfully:
\begin{Verbatim}[fontsize=\small]
lake env lean --run VerifyGeneralLog.lean
\end{Verbatim}
Its observed results were
\begin{Verbatim}[fontsize=\small]
Replaying 118179 endpoint dependencies in an empty kernel environment.
PASS: all 33 general-logarithm and counterexample endpoints.
PASS: their complete dependency closure rechecked
      in an empty kernel environment.
PASS: only propext, Classical.choice, and Quot.sound allowed as axioms.
\end{Verbatim}
The selected roots cover \(\log3\), rational and rational-power
logarithms, algebraic logarithms, both branches of the compatible
theorem, the arctangent and quadratic examples, and the real-input
counterexamples.
The checker collects dependencies of both types and values, follows
mutual declaration blocks, rejects unsafe or partial declarations
and nonstandard axioms, calls \texttt{Lean.Replay} in an environment
created by \texttt{mkEmptyEnvironment}, and checks that the selected
roots occur in the replayed environment.
This replay uses Lean's kernel in a fresh environment, not a second,
independently implemented proof checker.

\subsection{Pinned sources and reproduction}

The source package fixes the following revisions:
\begin{center}
\small
\begin{tabular}{@{}ll@{}}
\toprule
Component & Version or commit\\
\midrule
Lean & \texttt{4.34.1}\\
Lean commit & \texttt{5045d0056413266e57c625dcd7c365b10e377c52}\\
Mathlib & \texttt{d13f23b723b8a846827a245b89c10fc7d3f11612}\\
OpenAI mathematics & \texttt{adc7f1241b42e322a6451854ab7e4b4c146bf78a}\\
\bottomrule
\end{tabular}
\end{center}
The upstream tracked source was not modified.
The ancillary package contains the extension, selected upstream
source trees, \texttt{lean-toolchain}, \texttt{lakefile.toml},
the pinned \texttt{lake-manifest.json}, and the replay programs.
A SHA-256 inventory identifies the packaged files.
Only local source modules required by the logarithm formalization and
its verification targets are included.

After installing the specified Lean toolchain, enter the ancillary
Lean project directory and run
\begin{Verbatim}[fontsize=\small]
lake exe cache get
lake build LogarithmExtension.Verification
lake env lean --run VerifyGeneralLog.lean
\end{Verbatim}
Finish the build before starting the replay, and do not rebuild
imported modules while it is running.
The cache command obtains the Mathlib build cache; compiling its
sources is an alternative.
External toolchain and package downloads, adequate memory, and disk
space are required. The distributed archive contains sources rather
than a preinstalled Lean environment.

\subsection*{Acknowledgments}

The determinant method and much of its formal infrastructure are due
to the OpenAI preprint \cite{OAIpi}. Their reuse and adaptation are
credited throughout, and the upstream license notices are retained
in the source package.
The Lean and Mathlib communities supplied the proof assistant and
general mathematical libraries.
This work was developed with substantial assistance from OpenAI's
Codex, including exploration, implementation of the formal proofs,
and manuscript preparation.
